\documentclass[11pt,a4paper,fontset=fandol]{ctexart}

\usepackage[a4paper,top=2.25cm,bottom=2.45cm,left=2.25cm,right=2.25cm,headheight=18pt,headsep=0.55cm,footskip=24pt]{geometry}
\usepackage{amsmath,amssymb,amsthm}
\usepackage{fancyhdr}
\usepackage{lastpage}
\usepackage{enumitem}
\usepackage{titlesec}
\usepackage{xcolor}
\usepackage{hyperref}
\usepackage{bookmark}
\usepackage[most]{tcolorbox}
\usepackage{setspace}
\usepackage{array}

\hypersetup{
  colorlinks=true,
  linkcolor=blue!50!black,
  urlcolor=blue!50!black
}

\setstretch{1.13}
\setlength{\parindent}{2em}
\setlength{\parskip}{0.25em}
\allowdisplaybreaks
\setlist[enumerate]{itemsep=0.45em, topsep=0.35em, leftmargin=2.4em}
\setlist[itemize]{itemsep=0.25em, topsep=0.25em, leftmargin=1.9em}

\definecolor{mainblue}{RGB}{36,83,140}
\definecolor{lightblue}{RGB}{241,246,252}
\definecolor{lightgray}{RGB}{246,246,246}
\definecolor{rulegray}{RGB}{110,118,128}

\titleformat{\section}
  {\Large\bfseries\color{mainblue}}
  {\thesection.}
  {0.45em}
  {}
  [\vspace{0.2em}\color{rulegray}\titlerule]
\titlespacing*{\section}{0pt}{1.1em}{0.7em}

\pagestyle{fancy}
\fancyhf{}
\fancyhead[L]{\small 函数图像综合变换周测 B 卷}
\fancyhead[C]{\small 代数专题：综合判断与变式}
\fancyhead[R]{\small 刘文博}
\fancyfoot[C]{\small 第 \thepage 页 / 共 \pageref{LastPage} 页}
\renewcommand{\headrulewidth}{0.55pt}
\renewcommand{\footrulewidth}{0.35pt}

\newtcolorbox{infobox}[1]{
  breakable,
  colback=lightblue,
  colframe=mainblue,
  boxrule=0.7pt,
  arc=1mm,
  title=\textbf{#1},
  fonttitle=\bfseries
}

\newenvironment{solution}{\par\noindent\textbf{参考解答：}}{\hfill$\square$\par}
\newcommand{\answerspace}{\vspace{2.3cm}}
\newcommand{\biganswerspace}{\vspace{3.2cm}}

\begin{document}

\thispagestyle{empty}
\begin{center}
{\fontsize{22}{28}\selectfont\bfseries 函数图像综合变换周测 B 卷\par}
\vspace{0.4cm}
{\large 适合初中阶段函数专题提高训练\par}
\end{center}

\begin{infobox}{考试说明}
\begin{itemize}
  \item 测试时间：70--75 分钟
  \item 满分：100 分
  \item B 卷更强调顺序比较、分层解释和点的跟踪，请尽量写清楚“为什么这样变换”。
\end{itemize}
\end{infobox}

\section{试题}

\begin{enumerate}
  \item （12 分）请用“图像上的点跟着变化”的方法说明：为什么函数
  \[
  y=f(x)
  \]
  先向右平移 $a$ 个单位，再关于 \(y\) 轴对称后，可以写成
  \[
  y=f(-x-a).
  \]
  \biganswerspace

  \item （12 分）把函数
  \[
  y=x^2-2x
  \]
  先关于 \(y\) 轴对称，再向右平移 $3$ 个单位，最后向下平移 $2$ 个单位，求化简后的新方程。
  \answerspace

  \item （12 分）已知函数
  \[
  y=|x|
  \]
  经过综合变换后，顶点变成 \((-4,2)\)，图像开口向下。求新方程。
  \answerspace

  \item （12 分）已知点 \((2,5)\) 在函数
  \[
  y=f(x)
  \]
  的图像上。若新图像由原图像先关于 \(y\) 轴对称，再向右平移 $1$ 个单位得到，求对应点的新坐标。
  \answerspace

  \item （12 分）已知函数
  \[
  y=f(x)
  \]
  先向右平移 $2$ 个单位，再关于原点对称，写出所得新方程。
  \answerspace

  \item （12 分）比较下列两个结果是否一定相同，并说明理由：
  \begin{enumerate}
    \item 先关于原点对称，再向上平移 $3$ 个单位；
    \item 先向上平移 $3$ 个单位，再关于原点对称。
  \end{enumerate}
  \answerspace

  \item （14 分）将函数
  \[
  y=2x+1
  \]
  先关于 \(y\) 轴对称，再作一次竖直平移后，恰好经过原点。求第二次变化的方向、距离以及最终方程。
  \biganswerspace

  \item （14 分）已知原图像上的点 \((-1,2)\) 在函数
  \[
  y=f(x)
  \]
  上。若新图像由原图像先向左平移 $2$ 个单位，再关于 \(y\) 轴对称，最后关于 \(x\) 轴对称得到，求新图像的方程，并说明该点的对应变化。
  \biganswerspace
\end{enumerate}

\newpage
\section{详细解答}

\begin{enumerate}
  \item \begin{solution}
  设原图像上有一点 \((x_0,y_0)\)，因为它在函数
  \[
  y=f(x)
  \]
  上，所以有
  \[
  y_0=f(x_0).
  \]

  先向右平移 $a$ 个单位后，这个点变成
  \[
  (x_0+a,y_0).
  \]
  再关于 \(y\) 轴对称，变成
  \[
  (-(x_0+a),y_0)=(-x_0-a,y_0).
  \]

  若把新图像上的横坐标记作 $x$，那么对应原图像上的横坐标应为
  \[
  -x-a.
  \]
  因此新图像上的函数值是
  \[
  y=f(-x-a).
  \]

  所以所求方程为
  \[
  \boxed{y=f(-x-a)}.
  \]
  \end{solution}

  \item \begin{solution}
  先关于 \(y\) 轴对称：
  \[
  y=(-x)^2-2(-x)=x^2+2x.
  \]
  再向右平移 $3$ 个单位：
  \[
  y=(x-3)^2+2(x-3).
  \]
  最后向下平移 $2$ 个单位：
  \[
  y=(x-3)^2+2(x-3)-2.
  \]
  化简得
  \[
  y=x^2-4x+1.
  \]
  所以新方程为
  \[
  \boxed{y=x^2-4x+1}.
  \]
  \end{solution}

  \item \begin{solution}
  顶点从 \((0,0)\) 变成 \((-4,2)\)，说明向左平移 $4$ 个单位，再向上平移 $2$ 个单位；图像开口向下，说明关于 \(x\) 轴对称。

  所以新方程为
  \[
  \boxed{y=-|x+4|+2}.
  \]
  \end{solution}

  \item \begin{solution}
  点 \((2,5)\) 先关于 \(y\) 轴对称，变成 \((-2,5)\)。

  再向右平移 $1$ 个单位，变成 \((-1,5)\)。

  所以对应点的新坐标是
  \[
  \boxed{(-1,5)}.
  \]
  \end{solution}

  \item \begin{solution}
  先向右平移 $2$ 个单位，得到
  \[
  y=f(x-2).
  \]
  再关于原点对称，应写成
  \[
  y=-f(-(x)-2)=-f(-x-2).
  \]
  所以所求新方程为
  \[
  \boxed{y=-f(-x-2)}.
  \]
  \end{solution}

  \item \begin{solution}
  一般不相同。

  过程 1：
  \[
  y=f(x) \rightarrow y=-f(-x) \rightarrow y=-f(-x)+3.
  \]

  过程 2：
  \[
  y=f(x) \rightarrow y=f(x)+3 \rightarrow y=-(f(-x)+3)=-f(-x)-3.
  \]

  两个结果一般不同，因为关于原点对称既涉及横坐标变号，又涉及函数值变号，与竖直平移交换顺序后常数项会改变符号。
  \end{solution}

  \item \begin{solution}
  先关于 \(y\) 轴对称：
  \[
  y=2(-x)+1=-2x+1.
  \]
  设再进行一次竖直平移 $k$ 个单位，得到
  \[
  y=-2x+1+k.
  \]

  因为最终图像经过原点 \((0,0)\)，所以有
  \[
  0=1+k,
  \]
  解得
  \[
  k=-1.
  \]

  这说明第二次变化是向下平移 $1$ 个单位，最终方程为
  \[
  \boxed{y=-2x}.
  \]
  \end{solution}

  \item \begin{solution}
  按顺序改写方程：
  \[
  y=f(x) \rightarrow y=f(x+2) \rightarrow y=f(-x+2) \rightarrow y=-f(-x+2).
  \]
  所以新图像的方程为
  \[
  \boxed{y=-f(-x+2)}.
  \]

  再跟踪点的变化：
  \[
  (-1,2) \rightarrow (-3,2) \rightarrow (3,2) \rightarrow (3,-2).
  \]
  因此该点在新图像上对应为
  \[
  \boxed{(3,-2)}.
  \]
  \end{solution}
\end{enumerate}

\end{document}
